%=====================================================================
\chapter{Neural Network Foundations}\label{ch:neuralnets}
%=====================================================================
\locator{4}{Part IV --- Deep Learning and Modern AI}

\begin{outcomes}
By the end of this chapter you will be able to:
\begin{tight}
  \item \textbf{Compute} the output of a single neuron by hand.
  \item \textbf{Describe} the forward pass, loss computation, backpropagation and weight update as one training cycle.
  \item \textbf{Choose} an activation function and explain why non-linearity is essential.
  \item \textbf{Diagnose} a training run from its loss curves.
  \item \textbf{Decide} whether a problem calls for a neural network at all.
\end{tight}
\end{outcomes}

\begin{prereq}
\chref{math} (gradients, gradient descent, learning rate --- this chapter is that
machinery applied at scale), \chref{regression} (loss functions) and
\chref{mlfundamentals} (overfitting).
\end{prereq}

\begin{keyterms}
neuron $\bullet$ weight $\bullet$ bias $\bullet$ activation function $\bullet$
ReLU $\bullet$ sigmoid $\bullet$ softmax $\bullet$ hidden layer $\bullet$ forward
pass $\bullet$ loss $\bullet$ backpropagation $\bullet$ epoch $\bullet$ batch
$\bullet$ optimiser $\bullet$ vanishing gradient $\bullet$ dropout $\bullet$ early
stopping
\end{keyterms}

\section{One Neuron}

\begin{definitionbox}[Artificial Neuron]
A neuron computes a weighted sum of its inputs, adds a \term{bias}, and passes the
result through an \term{activation function}:
\[
z = \sum_{j=1}^{p} w_j x_j + b, \qquad a = f(z)
\]
The weights $w_j$ and bias $b$ are learned; the activation $f$ is chosen by you.
\end{definitionbox}

\dsfig{%
\begin{tikzpicture}[font=\scriptsize,
  inp/.style={circle, draw=secondary, fill=secondary!12, text=secondary!55!black,
              minimum size=0.75cm, font=\scriptsize, inner sep=0pt},
  arr/.style={-{Stealth[length=2mm]}, gray!65, line width=0.8pt}]

\node[inp] (x1) at (0,1.5)  {$x_1$};
\node[inp] (x2) at (0,0.5)  {$x_2$};
\node[inp] (x3) at (0,-0.5) {$x_3$};
\node[inp, draw=goldhl, fill=goldhl!14, text=goldhl!45!black] (b) at (0,-1.6) {$1$};

\node[rounded corners=4pt, draw=primary, fill=primary!10, text=primary!55!black,
      minimum width=1.9cm, minimum height=1.5cm, align=center, font=\scriptsize\bfseries]
      (sum) at (3.4,0.15) {$z=\sum w_jx_j$\\$+\,b$};

\node[rounded corners=4pt, draw=greenhl, fill=greenhl!12, text=greenhl!45!black,
      minimum width=1.7cm, minimum height=1.5cm, align=center, font=\scriptsize\bfseries]
      (act) at (6.4,0.15) {$f(z)$\\{\tiny activation}};

\node[circle, draw=accent, fill=accent!12, text=accent!70!black, minimum size=0.85cm,
      font=\scriptsize\bfseries] (out) at (8.9,0.15) {$a$};

\draw[arr] (x1) -- node[above, font=\tiny, pos=0.55, text=gray!55!black] {$w_1$} (sum);
\draw[arr] (x2) -- node[above, font=\tiny, pos=0.55, text=gray!55!black] {$w_2$} (sum);
\draw[arr] (x3) -- node[below, font=\tiny, pos=0.55, text=gray!55!black] {$w_3$} (sum);
\draw[arr, goldhl] (b) -- node[below, font=\tiny, pos=0.55, text=goldhl!45!black] {$b$} (sum);
\draw[arr] (sum) -- (act);
\draw[arr] (act) -- (out);

\node[font=\tiny\itshape, text=gray!55!black, align=center, text width=3.2cm] at (3.4,-2.1)
     {a linear model --- exactly the $z$ of \chref{classification}};
\node[font=\tiny\itshape, text=greenhl!45!black, align=center, text width=3.4cm] at (7.4,-2.1)
     {the non-linearity: without it, stacking layers gains nothing};
\node[font=\tiny\itshape, text=gray!55!black, anchor=west, text width=3.6cm] at (10.0,0.15)
     {feeds the next layer, or is the network's output};
\end{tikzpicture}}%
{A single neuron. The summation is a linear model; the activation is what makes depth
worthwhile. A neuron with a sigmoid activation \emph{is} logistic regression
(\chref{classification}).}{neuron}

\begin{worked}[One Neuron, by Hand]
Inputs $x = [0.8,\ 0.3,\ 0.5]$ (scaled attendance, quizzes, forum activity), weights
$w = [1.2,\ -0.5,\ 0.9]$, bias $b = -0.4$, activation ReLU.

\smallskip
\textbf{Weighted sum:}
\[
z = (1.2)(0.8) + (-0.5)(0.3) + (0.9)(0.5) - 0.4
  = 0.96 - 0.15 + 0.45 - 0.4 = 0.86
\]

\textbf{ReLU:} $a = \max(0,\ 0.86) = \mathbf{0.86}$. The neuron fires.

\smallskip
\textbf{Now change the input} to $x = [0.2,\ 0.9,\ 0.1]$:
\[
z = 0.24 - 0.45 + 0.09 - 0.4 = -0.52,
\qquad a = \max(0, -0.52) = \mathbf{0}
\]
The neuron is silent. This on/off behaviour is what lets a layer of neurons carve
the input space into regions --- and it is why ReLU produces \emph{sparse}
activations, with most neurons quiet for any given input.
\end{worked}

\section{Why Non-Linearity Is Not Optional}

\begin{alertbox}[The Collapse Argument]
Stack two linear layers: $y = W_2(W_1 x) = (W_2 W_1)x = W x$. The product of two
matrices is a matrix, so a hundred stacked linear layers compute exactly what one
linear layer computes. Without a non-linear activation between layers, depth buys
literally nothing. This single fact is why activation functions exist.
\end{alertbox}

\dsfig{%
\begin{tikzpicture}
\begin{groupplot}[
  group style={group size=4 by 1, horizontal sep=0.6cm},
  width=4.35cm, height=3.5cm,
  axis lines=middle, tick label style={font=\tiny},
  title style={font=\tiny\bfseries},
  xmin=-4, xmax=4, ymin=-1.4, ymax=1.9, xtick={-2,2}, ytick={-1,1},
]
\nextgroupplot[title={\color{greenhl}ReLU $\max(0,z)$}]
\addplot[greenhl, line width=1.3pt, domain=-4:0] {0};
\addplot[greenhl, line width=1.3pt, domain=0:1.8] {x};
\node[font=\tiny, text=gray!55!black, align=center, anchor=north] at (axis cs:-2.1,-0.25) {default\\for hidden layers};

\nextgroupplot[title={\color{secondary}Sigmoid $\sigma(z)$}]
\addplot[secondary, line width=1.3pt, domain=-4:4, samples=90] {1/(1+exp(-x))};
\node[font=\tiny, text=accent!80!black, align=center, anchor=north] at (axis cs:0,-0.35) {flat tails $\Rightarrow$\\vanishing gradient};

\nextgroupplot[title={\color{goldhl!60!black}tanh}]
\addplot[goldhl, line width=1.3pt, domain=-4:4, samples=90] {tanh(x)};
\node[font=\tiny, text=gray!55!black, align=center, anchor=north] at (axis cs:0,-0.35) {zero-centred\\version of sigmoid};

\nextgroupplot[title={\color{purplehl}Leaky ReLU}]
\addplot[purplehl, line width=1.3pt, domain=-4:0] {0.12*x};
\addplot[purplehl, line width=1.3pt, domain=0:1.8] {x};
\node[font=\tiny, text=gray!55!black, align=center, anchor=north] at (axis cs:-2.1,-0.3) {fixes\\``dead'' neurons};
\end{groupplot}
\end{tikzpicture}}%
{Four activation functions. ReLU is the default for hidden layers because it is cheap and
its gradient does not vanish for positive inputs; sigmoid and tanh saturate, which stalls
learning in deep networks.}{activations}

\rowcolors{2}{white}{lightbg}
\begin{longtable}{@{}L{2.6cm}L{4.4cm}L{7.2cm}@{}}
\toprule
\rowcolor{primary!15}
\textbf{Activation} & \textbf{Where} & \textbf{Why} \\
\midrule
\endhead
ReLU & Hidden layers (the default) & Cheap; gradient is 1 for $z>0$, so it does not vanish \\
Leaky ReLU & Hidden layers, if neurons die & Small negative slope keeps a gradient alive below zero \\
Sigmoid & \textbf{Output only}, binary classification & Produces a probability in $(0,1)$ \\
Softmax & \textbf{Output only}, multi-class & Produces probabilities across $k$ classes summing to 1 \\
None (linear) & \textbf{Output only}, regression & The prediction is an unbounded number \\
\bottomrule
\end{longtable}

\section{Layers and the Forward Pass}

\dsfig{%
\begin{tikzpicture}[font=\scriptsize,
  n/.style={circle, draw=#1, fill=#1!12, minimum size=0.52cm, inner sep=0pt},
  lbl/.style={font=\tiny\bfseries, align=center}]

\foreach \i in {1,...,4}{\node[n=secondary] (i\i) at (0,2.4-\i*0.85) {};}
\foreach \i in {1,...,5}{\node[n=primary]   (h\i) at (2.6,2.85-\i*0.85) {};}
\foreach \i in {1,...,5}{\node[n=primary]   (g\i) at (5.2,2.85-\i*0.85) {};}
\foreach \i in {1,...,2}{\node[n=accent]    (o\i) at (7.8,1.55-\i*0.85) {};}

\foreach \a in {1,...,4}{\foreach \b in {1,...,5}{
  \draw[gray!30, line width=0.35pt] (i\a) -- (h\b);}}
\foreach \a in {1,...,5}{\foreach \b in {1,...,5}{
  \draw[gray!30, line width=0.35pt] (h\a) -- (g\b);}}
\foreach \a in {1,...,5}{\foreach \b in {1,...,2}{
  \draw[gray!30, line width=0.35pt] (g\a) -- (o\b);}}

\node[lbl, text=secondary!55!black] at (0,2.15)  {INPUT\\{\tiny 4 features}};
\node[lbl, text=primary!55!black]   at (2.6,2.6) {HIDDEN 1\\{\tiny ReLU}};
\node[lbl, text=primary!55!black]   at (5.2,2.6) {HIDDEN 2\\{\tiny ReLU}};
\node[lbl, text=accent!70!black]    at (7.8,1.3) {OUTPUT\\{\tiny softmax}};

\draw[-{Stealth[length=2.6mm]}, greenhl, line width=1.6pt] (-0.9,-2.5) -- (8.7,-2.5);
\node[font=\tiny\bfseries, text=greenhl!45!black] at (3.9,-2.2) {FORWARD PASS: data flows right, producing a prediction};
\draw[-{Stealth[length=2.6mm]}, accent, line width=1.6pt] (8.7,-3.1) -- (-0.9,-3.1);
\node[font=\tiny\bfseries, text=accent!75!black] at (3.9,-3.42) {BACKPROPAGATION: error flows left, producing gradients};

\node[anchor=west, align=left, font=\tiny, text=gray!55!black, text width=4.2cm] at (9.1,0.2)
     {Parameters here:\\
      $4{\times}5 + 5 = 25$\\
      $5{\times}5 + 5 = 30$\\
      $5{\times}2 + 2 = 12$\\
      \textbf{67 numbers to learn}};
\end{tikzpicture}}%
{A small feed-forward network. Data moves right to produce a prediction; error moves left
to produce the gradient for every weight. Both directions run on every training
batch.}{network-layers}

\section{Training: The Four-Step Cycle}

\begin{definitionbox}[The Training Cycle]
\begin{tightnum}
  \item \textbf{Forward pass:} push a batch through the network to get predictions.
  \item \textbf{Loss:} measure how wrong they are (cross-entropy for classification,
        MSE for regression).
  \item \textbf{Backpropagation:} apply the chain rule backwards to get
        $\partial L / \partial w$ for every weight.
  \item \textbf{Update:} $w \leftarrow w - \alpha \nabla L$ --- exactly the
        gradient descent of \chref{math}.
\end{tightnum}
Repeat for every batch; one pass over the whole dataset is an \term{epoch}.
\end{definitionbox}

\begin{conceptbox}[Backpropagation in One Sentence]
Backpropagation is the chain rule from calculus, applied efficiently: it computes
how much each weight contributed to the final error by working backwards from the
output, reusing the partial results at each layer instead of recomputing them.
That efficiency is the entire reason deep networks are trainable at all.
\end{conceptbox}

\rowcolors{2}{white}{defbg}
\begin{longtable}{@{}L{3.4cm}L{4.4cm}L{6.4cm}@{}}
\toprule
\rowcolor{primary!15}
\textbf{Hyperparameter} & \textbf{Typical value} & \textbf{Effect if wrong} \\
\midrule
\endhead
Learning rate $\alpha$ & 0.001 (Adam) & \textbf{The one that matters most.} Too high diverges, too low crawls (\dsref{learning-rate}) \\
Batch size & 32--256 & Small = noisy but often generalises better; large = faster per epoch, needs more memory \\
Epochs & Until validation loss rises & Too few underfits; too many overfits --- use early stopping rather than guessing \\
Hidden layers & 1--3 for tabular data & More is not better on small data; deep networks need deep data \\
Neurons per layer & 16--256 & Too few underfits; too many overfits \\
Dropout & 0.2--0.5 & Randomly silences neurons during training, forcing redundancy \\
Optimiser & Adam & SGD needs more tuning; Adam is the safe default \\
\bottomrule
\end{longtable}

\section{Reading a Training Curve}

\dsfig{%
\begin{tikzpicture}
\begin{groupplot}[
  group style={group size=3 by 1, horizontal sep=0.95cm},
  width=5.35cm, height=4.2cm,
  axis lines=left, tick label style={font=\tiny}, label style={font=\tiny},
  title style={font=\scriptsize\bfseries}, xlabel={epoch}, ylabel={loss},
  xmin=0, xmax=50, ymin=0, ymax=2.2,
  legend style={font=\tiny, draw=gray!40, at={(0.98,0.98)}, anchor=north east},
]
\nextgroupplot[title={\color{greenhl}Healthy}]
\addplot[primary, line width=1.1pt, domain=1:50, samples=60] {1.9*exp(-0.075*x)+0.16};
\addplot[accent, line width=1.1pt, domain=1:50, samples=60] {1.9*exp(-0.070*x)+0.25};
\legend{train, validation}

\nextgroupplot[title={\color{accent}Overfitting}]
\addplot[primary, line width=1.1pt, domain=1:50, samples=60] {1.9*exp(-0.13*x)+0.03};
\addplot[accent, line width=1.1pt, domain=1:50, samples=60]
  {1.9*exp(-0.16*x)+0.34+0.00095*(x-14)^2};
\draw[greenhl, dashed, line width=1pt] (axis cs:14,0) -- (axis cs:14,2.0);
\node[font=\tiny, greenhl!50!black, anchor=south, align=center] at (axis cs:16,2.0) {stop here\\(early stopping)};

\nextgroupplot[title={\color{secondary}Learning rate too high}]
\addplot[primary, line width=1.1pt] coordinates {
 (1,1.9)(3,1.1)(5,1.7)(7,0.9)(9,1.8)(11,1.0)(13,1.9)(15,0.85)(17,2.0)(19,1.15)
 (21,1.85)(23,0.95)(25,1.95)(27,1.05)(29,1.9)(31,1.0)(33,2.05)(35,1.1)(37,1.9)
 (39,0.95)(41,2.0)(43,1.05)(45,1.85)(47,1.0)(49,1.95)};
\node[font=\tiny, secondary!60!black, align=center, anchor=north] at (axis cs:25,0.6)
     {loss bounces:\\lower $\alpha$};
\end{groupplot}
\end{tikzpicture}}%
{Three training runs. The gap between the two curves diagnoses overfitting; the shape of
the training curve alone diagnoses the learning rate. Plot both curves for every run you
do.}{training-curves}

\begin{pitfall}
\begin{tight}
  \item \textbf{Forgetting to scale inputs.} Neural networks are gradient-based, so
        unscaled features make the loss surface a narrow ravine and training stalls.
  \item \textbf{Sigmoid in hidden layers.} Its flat tails give near-zero gradients,
        so early layers stop learning --- the \term{vanishing gradient} problem.
        Use ReLU.
  \item \textbf{Training for a fixed number of epochs.} Use early stopping on
        validation loss.
  \item \textbf{Reaching for a network on 500 rows of tabular data.} Gradient
        boosting (\chref{features}) will beat it, train in seconds, and be
        explainable.
\end{tight}
\end{pitfall}

\begin{conceptbox}[When to Use a Neural Network --- and When Not To]
\textbf{Use one} for unstructured data --- images, audio, text, raw sequences ---
where features must be learned rather than engineered, and where you have tens of
thousands of examples or a pre-trained model to fine-tune.

\textbf{Do not} use one for ordinary tabular data of moderate size. On tabular
problems, gradient-boosted trees remain at least competitive and usually better,
while training faster and explaining themselves. This is one of the most reliably
reproduced results in applied machine learning, and ignoring it costs students a
great deal of time.
\end{conceptbox}

%---------------------------------------------------------------------
\begin{fourlenses}{Neural Networks}
\lensLA{Rarely the right tool for tabular engagement data --- a few thousand
students and 40 features is gradient-boosting territory. Networks earn their place
when the input is unstructured: assessing free-text reflections, or modelling the
raw clickstream \emph{sequence} rather than aggregates (\chref{timeseries}).
Whatever the accuracy, an unexplainable model driving interventions on students
runs into \chref{ethics}.}

\lensPM{Almost never appropriate. With 40 sprints and 67 parameters in even the toy
network above, the model has more parameters than observations. The honest answer
is a two-feature regression. Knowing when \emph{not} to use the powerful tool is
part of the competence being assessed.}

\lensIOT{Genuinely useful here, because raw sensor \emph{waveforms} are unstructured
and the fleet generates millions of examples. A small network can also be quantised
to run on-device (\chref{cloudedge}). Note the constraint that shapes everything:
the model must fit in kilobytes and infer in milliseconds, so architecture size is
dictated by the hardware, not by accuracy alone.}

\lensSEC{Used for malware classification from raw bytes and for sequence models over
session logs, where hand-engineered features are brittle against an adapting
adversary. But networks are also uniquely vulnerable: small, deliberately crafted
input perturbations can flip a confident prediction. That is adversarial machine
learning, and \chref{cybersecurity} treats it in full.}
\end{fourlenses}

%---------------------------------------------------------------------
\begin{lab}[A Network from Scratch, then with a Framework]
\begin{lstlisting}[language=Python]
import numpy as np

# ---- one neuron, no framework: verify the worked example --------------
x = np.array([0.8, 0.3, 0.5]); w = np.array([1.2, -0.5, 0.9]); b = -0.4
z = w @ x + b
print("z =", z, " relu(z) =", max(0.0, z))       # 0.86, 0.86

# ---- a tiny two-layer network trained by hand ------------------------
rng = np.random.default_rng(0)
X = rng.normal(size=(500, 3))
y = ((X[:, 0] + X[:, 2] - X[:, 1]) > 0).astype(float).reshape(-1, 1)

W1, b1 = rng.normal(size=(3, 8))*0.5, np.zeros((1, 8))
W2, b2 = rng.normal(size=(8, 1))*0.5, np.zeros((1, 1))
lr = 0.1

for epoch in range(400):
    h  = np.maximum(0, X @ W1 + b1)                 # forward: ReLU hidden
    p  = 1/(1 + np.exp(-(h @ W2 + b2)))             # forward: sigmoid out
    loss = -np.mean(y*np.log(p+1e-9) + (1-y)*np.log(1-p+1e-9))

    dz2 = (p - y) / len(X)                          # backprop
    dW2, db2 = h.T @ dz2, dz2.sum(0, keepdims=True)
    dh  = dz2 @ W2.T
    dz1 = dh * (h > 0)                              # ReLU derivative
    dW1, db1 = X.T @ dz1, dz1.sum(0, keepdims=True)

    W1 -= lr*dW1; b1 -= lr*db1; W2 -= lr*dW2; b2 -= lr*db2   # update
    if epoch % 100 == 0:
        print(f"epoch {epoch:3d}  loss {loss:.4f}")

# ---- the same thing with a framework, plus early stopping ------------
# from tensorflow import keras
# model = keras.Sequential([
#     keras.layers.Dense(32, activation="relu", input_shape=(3,)),
#     keras.layers.Dropout(0.3),
#     keras.layers.Dense(1,  activation="sigmoid")])
# model.compile(optimizer="adam", loss="binary_crossentropy", metrics=["AUC"])
# model.fit(X, y, epochs=200, validation_split=0.2,
#           callbacks=[keras.callbacks.EarlyStopping(patience=10,
#                                                    restore_best_weights=True)])
\end{lstlisting}

\textbf{Experiment:} set \texttt{lr = 5.0} and watch the loss bounce, reproducing
the right-hand panel of \dsref{training-curves}.
\end{lab}

%---------------------------------------------------------------------
\begin{checkpoint}
\begin{tightnum}
  \item A neuron has $w = [0.5, -1.0]$, $b = 0.3$, input $x = [2.0, 1.0]$, ReLU
        activation. Compute its output.
  \item Why does a network of purely linear layers gain nothing from depth?
  \item Training loss keeps falling while validation loss has been rising for 12
        epochs. What is happening and what should you do?
\end{tightnum}
\end{checkpoint}

%---------------------------------------------------------------------
\begin{chaptersummary}
\begin{tight}
  \item A neuron computes $f(\sum w_j x_j + b)$. With a sigmoid it is logistic
        regression; depth is what makes it more.
  \item Without non-linear activations, stacked layers collapse to a single linear
        layer.
  \item ReLU for hidden layers; sigmoid, softmax or linear at the output depending
        on the task.
  \item Training is forward pass $\rightarrow$ loss $\rightarrow$ backpropagation
        $\rightarrow$ update, repeated per batch.
  \item The learning rate is the hyperparameter that most often decides success.
        Use early stopping instead of a fixed epoch count.
  \item Scale inputs. And for moderate tabular data, prefer gradient boosting ---
        knowing when not to use a network is part of the skill.
\end{tight}
\end{chaptersummary}

\begin{reviewq}
  \item \textbf{(MCQ)} The standard activation for hidden layers is: (a)~sigmoid
        (b)~softmax (c)~ReLU (d)~linear.
  \item \textbf{(MCQ)} Backpropagation computes: (a)~the loss (b)~the gradient of
        the loss with respect to each weight (c)~the forward predictions
        (d)~the learning rate.
  \item \textbf{(Short)} Explain the vanishing gradient problem and why ReLU
        mitigates it.
  \item \textbf{(Short)} Define an epoch and a batch, and explain how they relate.
  \item \textbf{(Short)} Give two reasons to prefer gradient boosting over a neural
        network on a 5{,}000-row tabular dataset.
  \item \textbf{(Applied)} A network has 10 inputs, one hidden layer of 20 neurons,
        and 3 softmax outputs. Count the learnable parameters, showing your working.
  \item \textbf{(Applied)} Your training loss oscillates wildly and never settles.
        List three candidate causes in the order you would investigate them, with
        the check you would run for each.
\end{reviewq}
